% !TeX root = ../drafts/03-proving-primitives.tex
\section{Proving primitives}\label{sec:prelim}

\begin{definition}[Iverson bracket]\label{def:iverson}
For a proposition $P$, \ $[P]=1$ if $P$ holds and $[P]=0$ otherwise.
\end{definition}

\begin{definition}[Multilinear extension]\label{def:mle}
The \emph{multilinear extension} of $f:\cube{\kappa}\to\E$ is the unique $\mle{f}:\E^{\kappa}\to\E$ of degree at most one in each variable that agrees with $f$ on $\cube{\kappa}$.
\end{definition}


\begin{definition}[Equality indicator]\label{def:eq}
For $X,Y\in\E^{\kappa}$,
\[
  \eq(X,Y)= \prod_i\bigl(X_iY_i+(1+X_i)(1+Y_i)\bigr) = \prod_{i=1}^{\kappa}\bigl(1+X_i+Y_i\bigr).
\]
\end{definition}

$\eq$ is multilinear, and $\eq(X,Y)=[X=Y]$ on the boolean hypercube $\cube{2\kappa}$.

\begin{fact}[$\eq$ carries the interpolation weights]\label{fact:eq}
For an arbitrary point $r\in\E^{\kappa}$,
\[
  \mle{f}(r)=\sum_{x\in\cube{\kappa}}f(x)\,\eq(r,x).
\]
\end{fact}


Every committed column is in $\K$, while every random challenge is in $\E$, for soundness.

\begin{definition}[Table, column, virtual column]\label{def:column}
A \emph{column} of height $2^{\kappa}$ is a function $\cube{\kappa}\to\K$, and a \emph{table} is a group of columns of equal height. Each column is either \emph{committed} or \emph{virtual}: a virtual column is a fixed $\K$-linear combination of the committed columns of its own table. A $192$-bit \emph{value} (such as memory words) is indirectly committed by three columns, read as $V=V_0+V_1y+V_2y^{2}$. Addresses, registers and counters are single columns.
\end{definition}

\begin{fact}[Virtual columns are free]\label{fact:virtual}
Let $V=L(C)$ be virtual, $L$ a $\K$-linear map of the committed columns $C$ of its table, so $V$ is $\K$-valued like any other column. Extended $\E$-linearly, which is what evaluating at an $\E$-point asks of it, $L$ commutes with the multilinear extension: $\mle{V}(\zeta)=L\bigl(\mle{C}(\zeta)\bigr)$ for every $\zeta\in\E^{\kappa}$. One opening of the committed columns therefore evaluates every virtual column of the table, which is why the free increments $\gen\cdot\pc$ and $\gen\cdot\cnt$ cost nothing (\S\ref{sec:m3}).
\end{fact}

Two reductions turn statements about tables into evaluations of them.

\begin{fact}[Sumcheck~\cite{LFKN92}]\label{fact:sumcheck}
Over $\kappa$ rounds, sumcheck reduces a claim $\sum_{x\in\cube{\kappa}}P(x)=s$ on a polynomial $P$ in $\kappa$ variables to a single evaluation $P(\rho)$ at a random $\rho\in\E^{\kappa}$. A false claim survives with probability at most $\deg P\cdot\kappa/|\E|$.
\end{fact}

\begin{definition}[Zerocheck]\label{def:zerocheck}
To prove that $P$ vanishes on $\cube{\kappa}$, the verifier samples $r\in\E^{\kappa}$ and the two parties run sumcheck on
\[
  \sum_{x\in\cube{\kappa}}P(x)\,\eq(r,x)=0 .
\]
By Fact~\ref{fact:eq} that sum is the multilinear extension of $P|_{\cube{\kappa}}$ evaluated at $r$, so it vanishes identically in $r$ exactly when $P$ vanishes on the cube, and a $P$ that does not survives the random $r$ with probability at most $\kappa/|\E|$.
\end{definition}

\begin{definition}[Inner-product commitment scheme]\label{def:ipcs}
An \emph{inner-product commitment scheme} commits to a multilinear $g:\cube{M}\to\K$ and later proves, against any weight $W:\cube{M}\to\E$ whose extension the verifier can evaluate in $\mathrm{poly}(M)$ operations, a claim
\[
  \sum_{w\in\cube{M}}W(w)\,g(w)=c,
  \qquad c\in\E .
\]
The commitment is over $\K$, the opening over $\E$. A point evaluation $\mle{g}(r)=c$ is the special case $W=\eq(r,\cdot)$, by Fact~\ref{fact:eq}.
\end{definition}

WHIR~\cite{ACFY25} and BaseFold~\cite{ZCF23}/Ligerito~\cite{NA25} are such schemes. \S\ref{sec:rs-dense} fixes the one used here, and Annex~\ref{annex:pcs} specifies it with its round-by-round soundness, where Definition~\ref{def:ipcs} appears over a generic binary field as Definition~\ref{def:claim}.
